\section{Introduction}
\label{sec:intro}

Large language models are increasingly used as decision modules: a model reads a
case and emits an action that determines what a system does next. Under stochastic
decoding, repeated executions of the same input can produce different actions
\citep{Ouyang_2025,atil2025nondeterminismdeterministicllmsettings}. Operational
reliability therefore depends not only on whether a decision is correct, but whether
that decision recurs. Consider an accounts-payable system in which a model decides whether an invoice satisfies an approval policy. When the exact same invoice is labelled \texttt{approve} on one run and \texttt{reject} on the next, the decision is non-reproducible; repeatedly issuing the
same incorrect decision yields a reproducible but wrong system. In agentic systems, these failures can propagate through tool selection, workflow branching, and downstream context. Aggregate accuracy alone cannot distinguish these failure modes.

Existing approaches intervene at different parts of the stack. Batch-invariant kernels remove numerical sources of run-to-run variation
\citep{he2025nondeterminism}; fine-tuning changes learned behaviour
\citep{tian2023finetuninglanguagemodelsfactuality}; and constrained or
semantically guided decoding changes generation
\citep{loula2025syntactic,ahmed2025semanticprobabilisticcontrollanguage}.
Input-level methods instead alter the conditions presented to the model: prompt
search can shift output distributions
\citep{bhargava2024whatsmagicwordcontrol}, stability-guided refinement can improve
consistency \citep{chen2025stability}, and meaning-preserving changes in phrasing can alter both outputs and correctness
\citep{mizrahi-etal-2024-state,srikanth2024errorsnaturallanguagereasoning}. We take this further and ask whether interventions can move beyond statistical improvement to operational control: is there an intervention that does not merely make the correct action more likely, but establishes reproducible correctness across repeated single-pass executions? We answer this question by demonstrating that, for a fixed task, an explicit task specification can move failing model--input conditions to reproducible correctness without fine-tuning, retraining, decoder changes, retries, voting, or post-generation correction. These results position task specification as an engineering control surface: something that can be deliberately designed and tested against a declared
operational criterion, rather than treated as incidental wording around the model. Treating the specification, rather than only the solver, as an object of engineering
is standard outside machine learning. Requirements engineering
\citep{vanlamsweerde2009requirements}, formal specification
\citep{lamport2002specifying}, and contract-based design
\citep{meyer1992contract} treat requirements and interfaces as part of the
engineering problem. We adopt the same stance, treating task specification as an
engineering control surface for reproducible correctness rather than merely a means
of improving model behaviour incrementally.

Earlier work introduced \emph{interpretation drift} and showed that an explicit natural language task specification moved four frontier models from $30\%$ consistency (40 distinct hashes) to byte-identical
classifications (1 hash) across 40 runs
(Anonymous author, 2026). Here we extend that investigation to operational action selection using the car-wash prompt,\footnote{I need to wash my car, the car wash
is only 50 metres away, should I walk or drive there?} a public stress test used in multi-model comparisons and regularly highlighted as an example of jagged LLM capability. We define \texttt{DRIVE} as the ground-truth action and repeat the same task for ten independent single-pass executions per model--input condition. Of the 28 models measured at baseline, twenty-one are \emph{reproducibly incorrect}, producing \texttt{WALK} in all ten executions, and seven are \emph{non-reproducible}, producing both \texttt{DRIVE} and \texttt{WALK} across
the ten executions. Adding a natural-language task specification to the identical task moves 22 of these 28 failing model--input conditions to \emph{reproducible correctness}: 10/10 executions producing \texttt{DRIVE}. Conventional interventions produce substantially fewer such transitions and can also move non-reproducible models into reproducible incorrectness. A further condition states the correct action outright; it leaves seven models reproducibly incorrect, every one of which is reproducibly correct under the specification, so specifying the task and asserting the answer are not the same intervention. In a separate experiment on fourteen open-weight models, activation transplantation shows that whenever the specification reverses the final \texttt{DRIVE}--\texttt{WALK} preference, transferring the corresponding
specification-conditioned answer-position state into the untreated computation is sufficient to reverse that preference without the specification text being present
in the patched forward pass.

\paragraph{Contributions.}

\begin{itemize}\itemsep2pt

\item \textbf{From statistical improvement to engineered convergence.}
We show that input-level intervention need not be limited to improving average
performance. We adopt a criterion that admits no partial credit and sorts every
model--input condition into one of three operational states: \emph{reproducibly
correct} when all $R$ executions name the intended action, \emph{reproducibly
incorrect} when all $R$ name the same wrong one, and \emph{non-reproducible} when
the decision varies across them. The three are not points on a scale. A
reproducibly incorrect condition passes the cheapest check available at deployment
--- run the input twice and compare --- while being uniformly wrong, whereas a
non-reproducible one discloses its own unreliability to that same check. On the
car-wash task one shared specification moves $22$ of $28$ baseline-failing
conditions to reproducible correctness; the seven conventional prompt interventions
together move $2$, and five of the seven push at least one non-reproducible model
into reproducible incorrectness --- making its failure more consistent and harder to
detect at once.

\item \textbf{Substrate Engineering: natural language specification as a first-class engineering primitive.}
We introduce \emph{Substrate Engineering} as the practice of constructing and
testing natural-language task specifications against a declared reproducibility
criterion, treating the specification itself as a first-class engineering artifact
rather than incidental prompt wording. The resulting model-dependent successes and
failures expose the boundary of that control surface: the same specification can
establish reproducible correctness for one model--condition while remaining
insufficient for another. Reliability is therefore established for the
model--specification--task configuration, not attributed to the model in isolation.

\item \textbf{Under-specification as a measured quantity, not a verdict on the model.}
Reading a failure as evidence about the model presumes the input stated the task. We
make that assumption testable by turning the reliability requirement into a number.
Writing $p = \sigma(M)$ for the probability that a single draw names the intended
action, reproducible correctness over $R$ draws has probability $p^{R}$, so a
declared target fixes the margin an input must deliver: at $R=10$, $+2.63$ nats for
an even chance of $10/10$, $+4.55$ for $90\%$, and $+6.90$ for $99\%$. The distance
between a model's margin under a given input and that threshold is how much
specification is still missing --- in nats, for that model under that input. It is a
floor attached to the model--input pair, not to the model: none of the $28$ models
clears the requirement at baseline, $22$ clear it under the substrate, and among the
fourteen we can measure internally, eight exceed $+2.63$. The engineering question
becomes whether a given input specifies the task well enough for a chosen model to
meet a stated requirement, rather than whether that model is ``capable'' --- which
is not well posed until the input is named. We do not identify a minimal
specification, and the measured gap says how far there is to go, not what to add.

\item \textbf{A verified behavioural reference for internal causal analysis.}
We contribute a methodology that anchors internal causal analysis to an
experimentally verified behavioural transition rather than to an isolated successful
response. This fixes the behavioural object before asking what internal state is
causally responsible for changing it. In eight open-weight models, the substrate produces a verified transition from ten
incorrect executions to ten correct ones; in each, activation transplantation shows
that the conditioned answer-position state is sufficient, under the remaining
untreated computation, to reverse the measured action preference without the
specification text being present in the patched forward pass. A ninth model reverses that preference internally while its sampled behaviour remains unchanged---four of
ten executions name the intended action under both conditions. This dissociation shows why the behavioural reference matters: an internal preference reversal is not equivalent to reproducibly correct behaviour.
